% Fichier généré automatiquement par tools/generate_corrections.py.
% Source : RDM/RDM-05-Deformation/542_RdM/542_RdM.tex
% Statut : complet (2/2 question(s) renseignée(s), 0 reprise(s)).
\begin{corrigebox}[Corrigé extrait de la banque]
\CorrigeQuestion{1}
\textbf{Calcul du torseur de cohésion}

La poutre est composée de 2 tronçons :
\begin{itemize}
\item 1\ier tronçon : $\lambda\in[0,L]$
\item 2\ieme tronçon : $\lambda\in[L,2L]$
\end{itemize}

\textbf{1\ier tronçon}
\begin{itemize}
\item On isole la partie II.
\item BAME : 
\begin{itemize}
\item $\torseurstat{T}{I}{II}$
\item $\torseurstat{T}{p}{II}$
\end{itemize}
\item D'après le PFS, on a donc $\torseurstat{T}{I}{II} + \torseurstat{T}{p}{II} = \{0\}$ et donc 
$\torseurstat{T}{II}{I} =  \torseurstat{T}{p}{II}$.
\end{itemize}
On a donc $\forall \lambda\in[0,L]$, $\vectm{G}{F}{II}=  \left(3L/2-\lambda\right)\vect{x} \wedge - pL\vect{y} =- pL \left(\dfrac{3L}{2}-\lambda\right) \vect{z} $. 

Au final, $\torseurstat{T}{II}{I} = \torseurl{-pL\vect{y}}{- pL \left(\dfrac{3L}{2}-\lambda\right) \vect{z}}{G}$.

A SUIVRE...
\CorrigeQuestion{2}
\par\medskip\begin{center}\captionsetup{type=figure}
\centering

\end{center}\par\medskip
\end{corrigebox}
